\chapter*{Anexo B: Estructura completa del formato Open CVN JSON}
\addcontentsline{toc}{chapter}{Anexo B: Estructura completa del formato Open CVN JSON}
\label{cap:anexo_formato_json}

\setcounter{section}{0}
\renewcommand{\thesection}{B.\arabic{section}}
\setcounter{table}{0}
\renewcommand{\thetable}{B.\arabic{table}}

La Sección~\ref{sec:formato_json} introdujo la estructura raíz del formato
Open CVN JSON y su forma común para valores de referencia controlada. Este
anexo completa esa descripción con el detalle de cada bloque, tal como está
declarado en el artefacto canónico \texttt{schemas/open\_cvn.schema.json}
generado por el pipeline descrito en el
Capítulo~\ref{cap:pipeline}. El objetivo es ofrecer una referencia legible
por personas de la forma del formato, no sustituir el esquema JSON como
fuente de verdad para la validación automática.

\section{Bloque \texttt{metadata}}
\label{sec:anexo_formato_metadata}

El bloque \texttt{metadata} describe el documento, no su contenido
curricular, y contiene tres campos:

\begin{itemize}
  \item \texttt{language}: idioma del contenido curricular, expresado como
  código de idioma.

  \item \texttt{policy}: objeto con \texttt{name} y \texttt{version}, que
  identifica la política semántica aplicada durante la generación del
  documento, en el sentido descrito en la
  Sección~\ref{sec:pipeline_politica_semantica}.

  \item \texttt{source}: objeto con \texttt{format}, que identifica el
  origen del documento cuando procede de un proceso de importación.
\end{itemize}

\section{Bloque \texttt{curriculum}}
\label{sec:anexo_formato_curriculum}

El bloque \texttt{curriculum} organiza la información por área conceptual.
El campo \texttt{identity} es un único objeto que describe a la persona
titular del currículo; a diferencia del resto de campos, el esquema
canónico lo declara como un objeto abierto, sin una lista cerrada de
propiedades, porque la identidad agrupa datos de naturaleza heterogénea
(nombre, identificadores oficiales, sexo administrativo, trazas) que no
comparten una única forma repetida.

Los cinco campos restantes son listas de entradas repetidas:
\texttt{education}, \texttt{research}, \texttt{professional\_experience},
\texttt{achievements} y \texttt{other}. Cada entrada de estas listas
comparte una misma envoltura, resumida en la
Tabla~\ref{tab:anexo_formato_entrada}.

\begin{table}[htbp]
  \centering
  \small
  \renewcommand{\arraystretch}{1.1}
  \begin{tabular}{p{0.22\textwidth} p{0.68\textwidth}}
    \hline
    \textbf{Campo} & \textbf{Contenido} \\
    \hline
    \texttt{id} & Identificador local opcional de la entrada dentro del documento \\
    \texttt{type} & Identificador conceptual estable del tipo de entrada, independiente del código CVN de origen \\
    \texttt{data} & Objeto con los campos propios de la entrada, según su \texttt{type} \\
    \texttt{extensions} & Objeto opcional para metadatos específicos de una herramienta o proyecto concreto \\
    \texttt{trace} & Bloque opcional de trazabilidad hacia el paquete oficial CVN \\
    \hline
  \end{tabular}
  \caption{Envoltura común de las entradas de las listas curriculares.}
  \label{tab:anexo_formato_entrada}
\end{table}

Cuando está presente, el bloque \texttt{trace} puede incluir los códigos CVN
de origen (\texttt{cvn\_codes}), las rutas XML de las que procede
(\texttt{xml\_paths}), el nivel de confianza de la extracción
(\texttt{confidence}), y campos adicionales de apoyo a la trazabilidad
técnica: forma del dominio, tabla de referencia del manual funcional,
tipo de referencia semántica, patrón de serialización, y los artefactos y
ficheros de origen concretos. Este mismo mecanismo de trazabilidad es el
que permite, en la dirección inversa, relacionar cada campo del formato
Open CVN con el código CVN del que procede, tal como se describe en la
Sección~\ref{sec:pipeline_normalizacion}.

\section{Valores de referencia controlada}
\label{sec:anexo_formato_valores}

Los campos que remiten a una tabla de valores controlados comparten la
forma introducida en la Sección~\ref{sec:formato_json}: un objeto con
\texttt{code}, \texttt{label}, \texttt{source} y, cuando procede,
\texttt{raw\_value} o \texttt{uri}. El esquema canónico distingue dos
tratamientos posibles para esta forma, ya contrastados con evidencia
concreta en la Sección~\ref{sec:pipeline_politica_semantica}:

\begin{itemize}
  \item cuando la tabla de valores está cerrada y tiene evidencia
  suficiente, el campo \texttt{code} se emite como una enumeración cerrada
  en el esquema, como ocurre con \texttt{CVN\_SEX\_A};

  \item cuando la tabla no reúne evidencia suficiente para cerrarla sin
  perder valores válidos, el campo \texttt{code} permanece abierto, como
  ocurre con \texttt{CVN\_ENTITY\_TYPE}.
\end{itemize}

\section{Definiciones del esquema}
\label{sec:anexo_formato_definiciones}

El esquema canónico declara sus definiciones reutilizables bajo
\texttt{\$defs}. La Tabla~\ref{tab:anexo_formato_defs} resume su
distribución por área conceptual, verificada directamente sobre el
artefacto generado.

\begin{table}[htbp]
  \centering
  \small
  \renewcommand{\arraystretch}{1.1}
  \begin{tabular}{p{0.42\textwidth} p{0.15\textwidth} p{0.33\textwidth}}
    \hline
    \textbf{Área} & \textbf{Definiciones} & \textbf{Contenido} \\
    \hline
    \texttt{education} & 30 & formación académica e impartida \\
    \texttt{research} & 47 & actividad investigadora \\
    \texttt{professional\_experience} & 3 & experiencia profesional \\
    \texttt{achievements} & 2 & méritos y reconocimientos \\
    \texttt{other} & 20 & entradas sin área conceptual específica \\
    \texttt{identity} & 1 & identidad de la persona titular \\
    \texttt{core} & 1 & forma raíz del currículo \\
    vocabularios & 74 & valores de referencia controlada \\
    \hline
    \textbf{Total} & \textbf{182} & \\
    \hline
  \end{tabular}
  \caption{Distribución de las definiciones del esquema Open CVN JSON por área conceptual.}
  \label{tab:anexo_formato_defs}
\end{table}

Cada definición conserva, cuando existe evidencia, una traza no validante
hacia su origen mediante extensiones \texttt{x-open-cvn-*}: código CVN de
origen, rutas XML, forma del dominio conceptual, elegibilidad como
enumeración, referencia de origen y nivel de confianza. Estas extensiones
no participan en la validación JSON Schema en sentido estricto, pero
permiten auditar cada definición hasta su fuente oficial, en línea con el
principio de trazabilidad que atraviesa todo el trabajo.

Los artefactos citados en este anexo, el esquema canónico
\path{schemas/open_cvn.schema.json} y los documentos de ejemplo bajo
\path{examples/open_cvn/}, son reproducibles a partir del código fuente
del proyecto siguiendo el procedimiento descrito en el Anexo A.
